%%%%%%%%%%%%%%%%%%%%%%%%%%%%%%%%%
%
%       EXERCÍCIO
%
%%%%%%%%%%%%%%%%%%%%%%%%%%%%%%%%%

\ifdefstring{\atividade}{prova}{%
    \renewcommand{\valorquestao}{\ValorQ\ pontos}
}%

\begin{exercicioBanco}[\valorquestao]
Três empresas disputam a concessão para administrar o transporte público de uma cidade. A empresa A tem \(25\%\) de chance de vencer a licitação, a empresa B tem \(35\%\) e a empresa C tem \(40\%\). Caso vença a licitação, a probabilidade de investir em melhorias no sistema de transporte é de \(0{,}3\), \(0{,}5\) e \(0{,}8\), para as empresas A, B e C, respectivamente.

\begin{enumerate}[(a)]
    \item Qual é a probabilidade de não haver investimento em melhorias no transporte público?
    %
    \item Sabendo que houve investimento em melhorias, qual é a probabilidade de que a empresa \(A\) tenha vencido a licitação?
\end{enumerate}

\end{exercicioBanco}


